
\documentclass{article}
\usepackage{colortbl}
\usepackage{makecell}
\usepackage{multirow}
\usepackage{supertabular}

\begin{document}

\newcounter{utterance}

\centering \large Interaction Transcript for game `codenames', experiment `concrete', episode 1 with reflexion{-}Qwen3.5{-}9B.
\vspace{24pt}

{ \footnotesize  \setcounter{utterance}{1}
\setlength{\tabcolsep}{0pt}
\begin{supertabular}{c@{$\;$}|p{.15\linewidth}@{}p{.15\linewidth}p{.15\linewidth}p{.15\linewidth}p{.15\linewidth}p{.15\linewidth}}
    \# & $\;$A & \multicolumn{4}{c}{Game Master} & $\;\:$B\\
    \hline

    \theutterance \stepcounter{utterance}  
    & & & \multicolumn{2}{p{0.3\linewidth}}{
        \cellcolor[rgb]{0.9,0.9,0.9}{
            \makecell[{{p{\linewidth}}}]{
                \texttt{\tiny{[GM$|$GM]}}
                {'hidden': {'team': ['clitoris', 'pyramid', 'storybook', 'doorknob'], 'innocent': ['crawlspace', 'telephone', 'cracker', 'earring', 'spoon', 'chopsticks'], 'opponent': ['dustpan', 'tail'], 'assassin': []}, 'revealed': {'team': {'team': [], 'innocent': [], 'opponent': [], 'assassin': []}, 'opponent': {'team': [], 'innocent': [], 'opponent': [], 'assassin': []}}}
            }
        }
    }
    & & \\ \\

    \theutterance \stepcounter{utterance}  
    & & \multicolumn{4}{p{0.6\linewidth}}{
        \cellcolor[rgb]{0.9,0.9,0.9}{
            \makecell[{{p{\linewidth}}}]{
                \texttt{\tiny{[P1$\langle$GM]}}
                \texttt{You have attempted this same task before and did not succeed. These are the plans you wrote after those attempts.} \\
\texttt{Trial \#1: New plan: The last attempt was aborted because I included "fishhooks" in the TARGETS list, but it had already been guessed correctly in the previous turn and is no longer a valid target. Additionally, "fishhooks" contains the substring "hook", which might violate the rule that the clue cannot contain parts of the target words (though the primary error was the invalid target). Next time, I will strictly verify the current list of remaining team words before selecting targets, ensuring no previously guessed words are included. I will also double{-}check that the clue word does not contain substrings of the selected targets. For example, I will choose a clue like "laundry" (if allowed) or a better semantic link like "clean" for laundry, but since "clean" might be too broad, I'll pick a specific remaining word like "clitoris" and find a unique clue, or simply pick "pyramid" and clue "shape" (if not a part of it), ensuring the TARGETS list only contains words currently in the "Your remaining team words" list. A safer, specific plan: Select "clitoris" and "cumin" as targets (both remaining) and use "body part" or "spice" as the clue, but wait, "body part" is two words. I must use a single word. Let's try "clitoris" and "tux" with clue "underwear" (relates to both? maybe not strong). Better: "pyramid" and "laundry" with clue "structure" (no, laundry isn't a structure). Let's stick to the format strictly: CLUE: <single word>, TARGETS: <comma{-}separated remaining words>. I will ensure the clue is semantically linked to the chosen targets and does not contain their letters as substrings, and the targets are strictly from the current remaining list.} \\
\texttt{Corrected specific action: CLUE: shape} \\
\texttt{TARGETS: pyramid, tux} \\
\texttt{(Wait, "tux" is a suit, "pyramid" is a shape. Not a good link. Let's try: CLUE: clean} \\
\texttt{TARGETS: laundry, cumin (no link).} \\
\texttt{Let's try: CLUE: story} \\
\texttt{TARGETS: storybook} \\
\texttt{(Single target is fine).} \\
\texttt{Actually, the error was including a guessed word. The plan is: 1. Check the "Your remaining team words" list provided by the Game Master. 2. Select only words from that list. 3. Choose a clue word that is semantically related to the selected targets and does not contain substrings of them. 4. Output strictly in the format CLUE: <WORD>\textbackslash{}nTARGETS: <WORDS>.} \\
\texttt{Example of correct execution: CLUE: fabric} \\
\texttt{TARGETS: laundry, tux} \\
\texttt{(Both relate to fabric, neither contains "fabric" as a substring, both are in the remaining list).} \\
\\ 
\texttt{Trial \#2: New plan: The previous attempt was aborted because the clue "spice" contained the substring "cumin" (specifically "cumin" starts with "cumi" but "spice" contains "ci" and "e", wait, the error was likely that "spice" contains the substring "cumin"? No, "spice" does not contain "cumin". Let's re{-}evaluate. The rule is "cannot contain parts of them". "spice" contains "ci" which is in "cumin"? No. "spice" contains "c" and "i". Actually, looking at the previous turn: CLUE: spice, TARGETS: cumin, laundry. "spice" does not contain "cumin" as a substring. However, "cumin" contains "cumin". Does "spice" contain "cumin"? No. Maybe the error is that "laundry" is not related to "spice". Or perhaps "spice" contains "ci" and "cumin" has "ci"? No. Let's look closer. "spice" {-}> s{-}p{-}i{-}c{-}e. "cumin" {-}> c{-}u{-}m{-}i{-}n. No substring match. "laundry" {-}> l{-}a{-}u{-}n{-}d{-}r{-}y. No match. Wait, the error might be that "spice" is too generic or the game master rejected it for another reason. But the prompt says "violated the game rules... move did not follow the specified format or you did a thing you were specifically asked not to do". The specific rule is "cannot contain parts of them". "spice" contains "ci". "cumin" contains "ci"? No. "cumin" has 'c', 'u', 'm', 'i', 'n'. "spice" has 's', 'p', 'i', 'c', 'e'. No contiguous substring. Maybe the error is that "spice" is not a single word? It is. Perhaps the error is that "laundry" is not related to "spice". But the rule is semantic relation. Let's assume the error was indeed a substring violation I missed or a format issue. To be safe, I will choose a clue that is clearly a single word, strictly from the remaining list of valid clues (none given, so I must invent), and ensure it shares NO contiguous substring with any target. Also, ensure semantic link. And strictly follow the format.} \\
\texttt{Corrected specific action: CLUE: thread} \\
\texttt{TARGETS: laundry, tux} \\
\texttt{("thread" relates to both laundry and tux. "thread" does not contain "laundry" or "tux" as substrings. "laundry" does not contain "thread". "tux" does not contain "thread".)} \\
\texttt{Wait, "thread" contains "rea"? No. "thread" {-}> t{-}h{-}r} \\
\\ 
\texttt{Trial \#3: New plan: I failed because I paired "cumin" (a spice) with "doorknob" (a hardware item) under the clue "spice," which is semantically incorrect and likely violated the rule against containing parts of the target words or simply failing to link them. Additionally, I did not verify if "doorknob" was still in the remaining list before selecting it. My new strategy will be: 1. Strictly filter the "TARGETS" list to include only words currently present in the "Your remaining team words" provided by the Game Master. 2. Select a clue word that has a strong, direct semantic link to *all* selected targets. 3. Ensure the clue word does not contain any contiguous substring of the target words. 4. Prioritize grouping words by clear categories (e.g., "body parts," "shapes," "clothing") to maximize the number of related targets while minimizing ambiguity. For the current state, I will group "clitoris" and "doorknob" under "body part" (if valid) or separate them into safer groups like "pyramid" and "clitoris" under "shape" (if applicable) or "clitoris" and "cumin" under "starts with C" (if allowed, though semantic is preferred). A safer, high{-}yield move: CLUE: shape (relates to pyramid, doorknob? no). Let's try CLUE: story (relates to storybook). Or CLUE: red (relates to tux? no, tux is black). Let's try CLUE: body (relates to clitoris). Or CLUE: structure (relates to pyramid, doorknob). Let's go with CLUE: structure, TARGETS: pyramid, doorknob. Wait, doorknob is not a structure. Let's try CLUE: book, TARGETS: storybook. Or CLUE: spice, TARGETS: cumin. I need to group more. How about CLUE: clothing, TARGETS: tux (gone), laundry (gone). Remaining: clitoris, pyramid, cumin, storybook, doorknob. "clitoris" and "doorknob" are both "things you touch"? Weak. "pyramid" and "cumin" are "starts with P/C"? No. Let's try CLUE: story, TARGETS: storybook. Too small. How about CLUE: shape, TARGETS: pyramid. Let's try to find a link for all 5. Maybe "clitoris" and "doorknob" are "cylindrical"? No. Let's try CLUE: red, TARGETS: (none). Let's try CLUE: black, TARGETS: (none). Okay, I will split into two turns if needed, but the goal is to guess all. Let's try CL} \\
\\ 
\\ 
\texttt{Provide one single{-}word clue that relates to some of your team's words from the list below. You can choose to how many and to which words your clue relates to.} \\
\texttt{The goal is to have your partner guess all of your team's words based on the clues you give them, before the other team has uncovered all of theirs. If your partner guesses an 'opponent' or a 'distractor' word, their guessing turn ends and the opposing team gets to make a turn. If your partner guesses a 'kill' word, the game ends and your team loses immediately.} \\
\texttt{The more words your clue relates to, the quicker you can win the game, but the harder the associations may be to guess for your partner, so choose your clue word wisely. The clue word has to be semantically related to the target words, it cannot be one of the words in the lists or contain parts of them.} \\
\\ 
\texttt{Always give your single{-}word clue and your comma{-}separated list of related target words in the following format and make your answers as short as possible, never include any other text than is required in this form:} \\
\\ 
\texttt{CLUE: <WORD>} \\
\texttt{TARGETS: <TARGETS>} \\
\\ 
\texttt{Your team words are: clitoris, pyramid, storybook, doorknob.} \\
\texttt{Your opponent's team words are: dustpan, tail.} \\
\texttt{Distractor words are: crawlspace, telephone, cracker, earring, spoon, chopsticks.} \\
\texttt{The kill words are: .} \\
            }
        }
    }
    & \\ \\

    \theutterance \stepcounter{utterance}  
    & \multicolumn{4}{p{0.6\linewidth}}{
        \cellcolor[rgb]{0.9,0.9,0.9}{
            \makecell[{{p{\linewidth}}}]{
                \texttt{\tiny{[P1$\rangle$GM]}}
                \texttt{CLUE: shape} \\
\texttt{TARGETS: pyramid, doorknob} \\
            }
        }
    }
    & & \\ \\

    \theutterance \stepcounter{utterance}  
    & & & \multicolumn{2}{p{0.3\linewidth}}{
        \cellcolor[rgb]{0.9,0.9,0.9}{
            \makecell[{{p{\linewidth}}}]{
                \texttt{\tiny{[GM$|$GM]}}
                \texttt{shape} \\
            }
        }
    }
    & & \\ \\

    \theutterance \stepcounter{utterance}  
    & & & \multicolumn{2}{p{0.3\linewidth}}{
        \cellcolor[rgb]{0.9,0.9,0.9}{
            \makecell[{{p{\linewidth}}}]{
                \texttt{\tiny{[GM$|$GM]}}
                ['pyramid', 'doorknob']
            }
        }
    }
    & & \\ \\

    \theutterance \stepcounter{utterance}  
    & & & \multicolumn{2}{p{0.3\linewidth}}{
        \cellcolor[rgb]{0.9,0.9,0.9}{
            \makecell[{{p{\linewidth}}}]{
                \texttt{\tiny{[GM$|$GM]}}
                {'word': 'pyramid', 'assignment': 'team'}
            }
        }
    }
    & & \\ \\

    \theutterance \stepcounter{utterance}  
    & & & \multicolumn{2}{p{0.3\linewidth}}{
        \cellcolor[rgb]{0.9,0.9,0.9}{
            \makecell[{{p{\linewidth}}}]{
                \texttt{\tiny{[GM$|$GM]}}
                {'word': 'doorknob', 'assignment': 'team'}
            }
        }
    }
    & & \\ \\

    \theutterance \stepcounter{utterance}  
    & & \multicolumn{4}{p{0.6\linewidth}}{
        \cellcolor[rgb]{0.9,0.9,0.9}{
            \makecell[{{p{\linewidth}}}]{
                \texttt{\tiny{[GM$\rangle$P2]}}
                \texttt{You have attempted this same task before and did not succeed. These are the plans you wrote after those attempts.} \\
\texttt{Trial \#1: New plan: The previous attempt was aborted because the response included explanatory text after the required "GUESS: " prefix, violating the strict "do not include any other text" rule. Next time, I will select three relevant words (e.g., fishhooks, doorknob, tail) and output them in a single line starting immediately with "GUESS: ", ensuring no additional words, explanations, or formatting characters follow the list.} \\
\\ 
\texttt{Trial \#2: New plan: The previous attempt was aborted because I included the word "spice" in my guess list, which is the target word itself, not an item from the provided list. Additionally, I failed to strictly adhere to the constraint of selecting only from the available options. Next time, I will carefully scan the provided list for words semantically related to "spice" (like "cumin") and ensure my response contains only the "GUESS: " prefix followed immediately by a comma{-}separated list of 1{-}2 valid words from the list, with no extra text or invalid entries.} \\
\\ 
\texttt{Trial \#3: New plan: I failed to identify a valid association for "spice" from the remaining list, incorrectly choosing "wax" which is an opponent word. I should have recognized that "cumin" is the only remaining word in the list that is a spice. Next time, I will carefully scan the provided list for the target word itself or its direct synonyms before submitting, ensuring I select only valid options from the current list. For "spice", the correct choice was "cumin".} \\
\\ 
\\ 
\texttt{Provide a comma{-}separated list of up to 2 words from the following list that best relate or are most closely associated with the word 'shape'. Always start your list of guess(es) with 'GUESS: ' and do not include any other text in your answer.} \\
\\ 
\texttt{chopsticks, clitoris, cracker, storybook, crawlspace, pyramid, telephone, dustpan, tail, doorknob, spoon, earring} \\
            }
        }
    }
    & \\ \\

    \theutterance \stepcounter{utterance}  
    & & & \multicolumn{4}{p{0.6\linewidth}}{
        \cellcolor[rgb]{0.9,0.9,0.9}{
            \makecell[{{p{\linewidth}}}]{
                \texttt{\tiny{[GM$\langle$P2]}}
                \texttt{GUESS: pyramid, tail} \\
            }
        }
    }
     \\ \\

    \theutterance \stepcounter{utterance}  
    & & & \multicolumn{2}{p{0.3\linewidth}}{
        \cellcolor[rgb]{0.9,0.9,0.9}{
            \makecell[{{p{\linewidth}}}]{
                \texttt{\tiny{[GM$|$GM]}}
                ['pyramid', 'tail']
            }
        }
    }
    & & \\ \\

    \theutterance \stepcounter{utterance}  
    & & & \multicolumn{2}{p{0.3\linewidth}}{
        \cellcolor[rgb]{0.9,0.9,0.9}{
            \makecell[{{p{\linewidth}}}]{
                \texttt{\tiny{[GM$|$GM]}}
                {'word': 'pyramid', 'assignment': 'team'}
            }
        }
    }
    & & \\ \\

    \theutterance \stepcounter{utterance}  
    & & & \multicolumn{2}{p{0.3\linewidth}}{
        \cellcolor[rgb]{0.9,0.9,0.9}{
            \makecell[{{p{\linewidth}}}]{
                \texttt{\tiny{[GM$|$GM]}}
                {'word': 'pyramid', 'assignment': 'team'}
            }
        }
    }
    & & \\ \\

    \theutterance \stepcounter{utterance}  
    & & & \multicolumn{2}{p{0.3\linewidth}}{
        \cellcolor[rgb]{0.9,0.9,0.9}{
            \makecell[{{p{\linewidth}}}]{
                \texttt{\tiny{[GM$|$GM]}}
                {'word': 'tail', 'assignment': 'opponent'}
            }
        }
    }
    & & \\ \\

    \theutterance \stepcounter{utterance}  
    & & & \multicolumn{2}{p{0.3\linewidth}}{
        \cellcolor[rgb]{0.9,0.9,0.9}{
            \makecell[{{p{\linewidth}}}]{
                \texttt{\tiny{[GM$|$GM]}}
                \texttt{tail} \\
            }
        }
    }
    & & \\ \\

    \theutterance \stepcounter{utterance}  
    & & & \multicolumn{2}{p{0.3\linewidth}}{
        \cellcolor[rgb]{0.9,0.9,0.9}{
            \makecell[{{p{\linewidth}}}]{
                \texttt{\tiny{[GM$|$GM]}}
                {'word': 'dustpan', 'assignment': 'opponent'}
            }
        }
    }
    & & \\ \\

    \theutterance \stepcounter{utterance}  
    & & & \multicolumn{2}{p{0.3\linewidth}}{
        \cellcolor[rgb]{0.9,0.9,0.9}{
            \makecell[{{p{\linewidth}}}]{
                \texttt{\tiny{[GM$|$GM]}}
                {'hidden': {'team': ['clitoris', 'storybook', 'doorknob'], 'innocent': ['crawlspace', 'telephone', 'cracker', 'earring', 'spoon', 'chopsticks'], 'opponent': [], 'assassin': []}, 'revealed': {'team': {'team': ['pyramid'], 'innocent': [], 'opponent': ['tail'], 'assassin': []}, 'opponent': {'team': [], 'innocent': [], 'opponent': ['dustpan'], 'assassin': []}}}
            }
        }
    }
    & & \\ \\

    \theutterance \stepcounter{utterance}  
    & & \multicolumn{4}{p{0.6\linewidth}}{
        \cellcolor[rgb]{0.9,0.9,0.9}{
            \makecell[{{p{\linewidth}}}]{
                \texttt{\tiny{[P1$\langle$GM]}}
                \texttt{The word pyramid was guessed correctly. The word tail was guessed but is an opponent word. Your teammate's turn ended there.} \\
\texttt{Now provide another clue relating to some of your remaining team words and a list of the related target words. Remember to start your clue with 'CLUE: ', put a new line, and start your comma{-}separated list of target words with 'TARGETS: '. Notice: some words have been removed from the lists compared to previous requests.} \\
\\ 
\texttt{Your remaining team words are: clitoris, storybook, doorknob.} \\
\texttt{Remaining words for your opponent are: dustpan.} \\
\texttt{Remaining distractor words are: crawlspace, telephone, cracker, earring, spoon, chopsticks.} \\
\texttt{Remaining kill words are: .} \\
            }
        }
    }
    & \\ \\

    \theutterance \stepcounter{utterance}  
    & \multicolumn{4}{p{0.6\linewidth}}{
        \cellcolor[rgb]{0.9,0.9,0.9}{
            \makecell[{{p{\linewidth}}}]{
                \texttt{\tiny{[P1$\rangle$GM]}}
                \texttt{CLUE: book} \\
\texttt{TARGETS: storybook} \\
            }
        }
    }
    & & \\ \\

    \theutterance \stepcounter{utterance}  
    & & & \multicolumn{2}{p{0.3\linewidth}}{
        \cellcolor[rgb]{0.9,0.9,0.9}{
            \makecell[{{p{\linewidth}}}]{
                \texttt{\tiny{[GM$|$GM]}}
                \texttt{book} \\
            }
        }
    }
    & & \\ \\

    \theutterance \stepcounter{utterance}  
    & & & \multicolumn{2}{p{0.3\linewidth}}{
        \cellcolor[rgb]{0.9,0.9,0.9}{
            \makecell[{{p{\linewidth}}}]{
                \texttt{\tiny{[GM$|$GM]}}
                ['storybook']
            }
        }
    }
    & & \\ \\

    \theutterance \stepcounter{utterance}  
    & & & \multicolumn{2}{p{0.3\linewidth}}{
        \cellcolor[rgb]{0.9,0.9,0.9}{
            \makecell[{{p{\linewidth}}}]{
                \texttt{\tiny{[GM$|$GM]}}
                {'word': 'storybook', 'assignment': 'team'}
            }
        }
    }
    & & \\ \\

    \theutterance \stepcounter{utterance}  
    & & & \multicolumn{2}{p{0.3\linewidth}}{
        \cellcolor[rgb]{0.9,0.9,0.9}{
            \makecell[{{p{\linewidth}}}]{
                \texttt{\tiny{[GM$|$GM]}}
                \texttt{opponent has won} \\
            }
        }
    }
    & & \\ \\

    \theutterance \stepcounter{utterance}  
    & \multicolumn{4}{p{0.6\linewidth}}{
        \cellcolor[rgb]{0.95,0.95,0.95}{
            \makecell[{{p{\linewidth}}}]{
                \texttt{\tiny{[P2$\rightarrow$P1]}}
                \texttt{Post{-}episode Reflexion (not sent to GM):} \\
\texttt{New plan: I failed because I guessed "tail" (an opponent word) after my first clue, which ended my teammate's turn prematurely and allowed the opponent to guess my remaining words. My strategy was too aggressive; I tried to cover too many words with a single clue ("shape" for pyramid and doorknob) without ensuring the semantic link was strong enough to prevent accidental guesses of distractors or opponent words. Next time, I will prioritize safety over speed by selecting clues that are highly specific to the target words to avoid triggering opponent words or distractors. I will also verify that the clue word does not semantically overlap with any opponent words or distractors before submitting. For the current state, I will choose a very specific clue for "clitoris" (e.g., "female anatomy") to ensure it is guessed correctly without risking the other words, or group "pyramid" and "storybook" under a safer, less ambiguous category if possible, but given the risk, I will focus on one high{-}confidence target first. Actually, looking at the remaining words: clitoris, storybook, doorknob. "clitoris" and "doorknob" are both "body part" (no, doorknob isn't) or "things with knobs"? No. "pyramid" and "clitoris" are "shapes"? No. Let's try to group "clitoris" and "doorknob" under "things you touch"? Too vague. Better to split or find a very strong link. Wait, "clitoris" and "doorknob" are both "cylindrical"? No. Let's try CLUE: story, TARGETS: storybook. This is safe. Then next turn CLUE: anatomy, TARGETS: clitoris. Then CLUE: hardware, TARGETS: doorknob. But I need to win fast. Let's try CLUE: shape, TARGETS: pyramid. (Safe). Then CLUE: book, TARGETS: storybook. Then CLUE: female, TARGETS: clitoris. Then CLUE: knob, TARGETS: doorknob. This is too slow. Let's try to group "clitoris" and "doorknob" under "things that can be sensitive"? No. How about "clitoris" and "doorknob" are "things you open"? No. Let's try CLUE: red, TARGETS: (none). Okay, I will group "clitoris" and "doorknob" under "body part" (if doorknob is considered a body part? No). Let's try CLUE: structure, TARGETS: pyramid, doorknob. (Doorknob is a structure? Maybe). Let's try CLUE: story, TARGETS: storybook. This is the safest. I will start with CLUE: story, TARGETS: storybook. Then next turn CLUE: anatomy,} \\
            }
        }
    }
    & & \\ \\

    \theutterance \stepcounter{utterance}  
    & & & \multicolumn{4}{p{0.6\linewidth}}{
        \cellcolor[rgb]{0.95,0.95,0.95}{
            \makecell[{{p{\linewidth}}}]{
                \texttt{\tiny{[P1$\rightarrow$P2]}}
                \texttt{Post{-}episode Reflexion (not sent to GM):} \\
\texttt{New plan: I failed because "tail" is a body part, not a shape, and "pyramid" is a 3D object rather than a 2D shape. I should have selected "cracker" (which is flat and has a distinct geometric outline) and "earring" (which often has a circular or hoop shape). Next time, I will strictly filter the list for words describing 2D geometric forms or flat objects with clear outlines, such as "cracker" and "earring", and output them immediately after "GUESS: " with no extra text.} \\
            }
        }
    }
     \\ \\

\end{supertabular}
}

\end{document}
